\section*{Abstract}
This report describes a methodology for estimating aerodynamic drag area ($C_DA$), downforce area ($C_LA$), and rolling resistance coefficient ($C_{rr}$) from publicly available F1 telemetry using the FastF1 library.

Coast-down segments from free-practice sessions (FP1--FP3) are fitted to an ODE model incorporating MGU-K energy recovery. $C_LA$ is estimated independently via GPS-derived lateral acceleration from qualifying push laps. Constructor and circuit selection is data-driven: all 20 drivers across all 24 rounds of the 2024 calendar are surveyed on two joint metrics (coast-down segment count and qualifying-lap $C_LA$ sample count), and the constructor with the highest aggregate score is selected. Aston Martin (driver~14, Alonso, AMR24) ranks first; the five circuits with the best joint data quality spanning the full downforce range are chosen: Jeddah, Spa-Francorchamps, Silverstone, Yas Marina, and Suzuka.

A five-point aero polar ($C_LA$ vs.\ $C_DA$) fitted by OLS yields a slope of $3.0$ (90\% CI: $-2.4$ to $+5.8$~m$^2$/m$^2$; 1.09$\sigma$ from zero). The slope is positive---consistent with the expected physical relationship between wing drag and downforce---but the confidence interval includes zero, so a definitive efficiency figure cannot be stated. The dominant uncertainty source is the assumed tyre friction coefficient $\mu = 1.8$ in the $C_LA$ estimator ($\pm0.4$--$0.7$~m$^2$ systematic per circuit). DRS drag reduction is investigated via race trap speeds but found to be inseparable from slipstream effects in the available data.

An investigation into 2026 active aero measurements reveals that the F1 timing API does not broadcast active aero state, blocking mode-split analysis.
